En base al análisis realizado en el capítulo anterior, se establecen las metas fundamentales de este Trabajo de Fin de Grado. El propósito principal de este proyecto es dar respuesta a la sobresaturación y fragmentación del mercado mediante la creación de una obra que actúe como puente entre diferentes perfiles de usuario y así impulsar además la llegada de nuevos jugadores. 

Por tanto, el \textbf{objetivo general} de este proyecto es diseñar y desarrollar un prototipo jugable (\textit{Vertical Slice}) de \textbf{Reset}, un videojuego en 2.5D con estética \textit{Pixel Art} para PC que unifica múltiples paradigmas lúdicos bajo un mismo hilo conductor, sirviendo como demostración técnica y de diseño de las competencias adquiridas durante el grado.

Para alcanzar esta meta global, se definen los siguientes \textbf{objetivos específicos}:

\begin{itemize}
    \item \textbf{Analizar y establecer los pilares de un diseño estético coherente:} Definir y documentar de forma razonada la dirección artística y la usabilidad de las interfaces de usuario (UI/UX) basadas en el estilo Pixel Art, garantizando una sólida cohesión visual que justifique la coexistencia de mundos y géneros lúdicos dispares bajo una misma identidad estética.
    
    \item \textbf{Desarrollar una narrativa cohesiva y contextualizada:} Crear un marco narrativo meta-referencial y un sistema de diálogos interactivos que justifiquen y conecten orgánicamente los cambios de géneros y mecánicas entre mundos, enriqueciendo la inmersión del jugador a través de referencias y homenajes a obras icónicas de la industria.
    
    \item \textbf{Implementar un diseño de niveles estructurado en capas:} Aplicar la teoría de la fluidez para balancear la dificultad, construyendo una ruta principal accesible con su debida curva de aprendizaje para jugadores noveles y diseñando retos secundarios de alta exigencia para retener el interés de los jugadores veteranos.

    \item \textbf{Crear un movimiento y control en 2D adaptado al género de plataformas:} Diseñar y calibrar las mecánicas de movimiento dinámico del jugador aprovechando el motor de físicas, integrando algoritmos de \textit{Game Feel} para asegurar una experiencia reactiva, precisa y satisfactoria.

    \item \textbf{Desarrollar una arquitectura escalable de Inteligencia Artificial:} Estructurar los agentes enemigos mediante principios de programación orientada a objetos (herencia y polimorfismo) e implementar algoritmos de navegación y de toma de decisiones reactiva para garantizar comportamientos orgánicos y adaptativos.

    \item \textbf{Implementar un sistema robusto de persistencia de datos y gestión de estado:} Crear una arquitectura de guardado y carga centralizada que serialice en formato JSON el estado de todas las variables, favoreciendo la continuidad de la experiencia de forma compacta y segura.

   \item \textbf{Integrar sistemas mecánicos específicos para cada género:} Diseñar e implementar las lógicas particulares de cada entorno para conformar una experiencia de juego sólida, inmersiva y coherente en su conjunto.
\end{itemize}